\section{Codifica lossless}
\subsection{Introduzione alle codifiche lossless}
La codifica lossless consiste nel rappresentare una sequenza di simboli emessi da una sorgente tramite una stringa di bit
in modo perfettamente invertibile, cioè senza perdita di informazione. L'obiettivo delle codifiche lossless è quello di
rappresentare i simboli nel minor numero di bit possibile (preservando l'informazione) in modo da ridurre la quantità di
dati da trasmettere o memorizzare.

\subsubsection*{Definizioni di base}
\[\begin{array}{c c}
	\displaystyle \mathcal{X} = \left\{ x_1, x_2, \dots x_M \right\}
	& \quad \begin{array}{c}
		\text{alfabeto con } M \text{ simboli emessi da una sorgente } X, \\
		\text{ciascuno con la sua probabilità } p_i = P(X=x_i)
	\end{array}
	\\[12pt]

	\displaystyle C : x_i \in \mathcal{X} \to c_i \in \{0,1\}^*
	& \quad \begin{array}{c}
		\text{codice } C \text{ definito come funzione di codifica che associa} \\
		\text{una stringa di bit (o codeword)} c_i \text{ a ciascun simbolo } x_i \\
		\text{ogni codeword ha la sua lunghezza } l_i = |c_i|
	\end{array}
	\\[17pt]

	\displaystyle \mathcal{L} = \sum_{i=1}^{M} p_i l_i
	& \quad \begin{array}{c}
		\text{lunghezza media del codice calcolata come media delle lunghezze} \\
		\text{di ciascuna codeword pesate per le relative probabilità dei simboli}
	\end{array}
\end{array}\]

\subsubsection*{Codici univocamente decodificabili}
Per garantire assenza di perdita di informazione, il codice deve essere univocamente decodificabile, cioè ogni sequenza di
codewords senza separatori deve essere decodificabile in un'unica sequenza di simboli. Di seguito alcuni esempi di codici
univocamente decodificabili:
\begin{itemize}
	\item \textbf{fixed length coding - FLC}: ogni simbolo è associato ad una codeword di lunghezza fissa, la decodifica è
	estremamente facile, ma non è ottimale in quanto non tiene conto della distribuzione di probabilità dei simboli
	\item \textbf{prefix coding}: nessuna codeword è prefisso di un'altra codeword, quindi ogni volta che si legge una codeword,
	la si può decodificare immediatamente essendo certi che la lettura dei successivi bit non porterà a una codeword più lunga,
	permette di ottenere codici ottimi se si tiene conto della distribuzione di probabilità dei simboli
	\item \textbf{unary coding}: codice in cui le codewords sono sequenze di 1 (o 0) delimitate all'inizio o alla fine da un bit 0
	(o 1), la decodifica consiste nel contare il numero di bit uguali della sequenza prima di incontrare il bit delimitatore,
	che indica la fine della codeword o l'inizio della successiva
\end{itemize}

\subsubsection*{Entropia di una variabile aleatoria}
Data una variabile aleatoria \(X\) che assume valori discreti \(x_i\) di un alfabeto \(\mathcal{X}\), ciascuno con probabilità
\(p_i = P(X=x_i)\), si definiscono:
\[\begin{array}{c c}
	\displaystyle I(x_i) = \log_2 \frac{1}{p_i} = - \log_2 p_i
	& \quad \begin{array}{c}
		\text{informazione associata all'evento } X=x_i
	\end{array}
	\\[12pt]
	\displaystyle H(X) = E[I(X)] = \sum_{i=1}^{M} p_i \, I(x_i) = \sum_{i=1}^{M} p_i \log_2 \frac{1}{p_i}
	& \quad \begin{array}{c}
		\text{entropia della variabile aleatoria } X
	\end{array}
\end{array}\]
L'entropia rappresenta la misura della quantità di informazione media contenuta in un simbolo emesso dalla sorgente e
può essere interpretata come il grado di incertezza associato alla variabile aleatoria \(X\). L'entropia è massima quando
gli \(M\) simboli sono equiprobabili, con \(p_i = 1/M\), e vale \(H(X) = \log_2 M\). Viceversa quando si hanno pochi simboli
molto probabili e molti simboli poco probabili (distribuzioni sparse), l'entropia è molto bassa.

\subsubsection*{Teorema di Shannon sulla codifica di sorgente}
La lunghezza media \(\mathcal{L}^*\) di un codice decodificabile ottimo per una sorgente \(X\) è vincolata dalla relazione:
\[H(X) \; \leq \; \mathcal{L}^* \, < \; H(X) + 1 \qquad\quad \text{con } \mathcal{L}^* = H(X) \; \Longleftrightarrow \; p_i = 2^{-l_i} \text{ (per distribuzioni diadiche)}\]

\subsubsection*{Codifica a blocchi}
Per sfruttare la correlazione tra simboli consecutivi emessi dalla sorgente, è possibile raggruppare \(K\) simboli
consecutivi in un blocco e codificare ogni blocco come un simbolo unico.

Analizzando la formula del teorema di Shannon, si osserva che la lunghezza media per simbolo \(\mathcal{L}_S^*\) del codice
tende a convergere verso il limite teorico dettato dall'entropia all'aumentare di \(K\), siccome il bit di overhead \(+ 1\)
viene distribuito sui \(K\) simboli di ogni blocco.
\[H(X^K) \leq L^* < H(X^K) + 1 \;\; \implies \;\; \frac{H(X^K)}{K} \leq \mathcal{L}_S^* < \frac{H(X^K)}{K} + \frac{1}{K} \quad \text{con} \; \mathcal{L}_S^* = \frac{\mathcal{L}^*}{K}\]
\[\lim_{K \to \infty} \frac{H(X^K)}{K} = \lim_{K \to \infty} \frac{H(X^K)}{K} + \frac{1}{K} = \lim_{K \to \infty} \mathcal{L}_S^*\]

\subsubsection*{Tasso entropico come limite assoluto delle codifiche lossless}
Si osserva inoltre che l'entropia per blocco \(H(X^K)\) decresce all'aumentare della lunghezza del blocco \(K\) o al più
rimane costante se i simboli presi singolarmente sono indipendenti. Di conseguenza anche l'entropia per simbolo \(H(X^K)/K\)
e la relativa lunghezza media per simbolo \(\mathcal{L}_S^*\) decrescono fino a raggiungere il limite teorico assoluto per
le codifiche lossless detto tasso entropico \(\mathcal{H}(X)\) della sorgente \(X\).
\[\mathcal{H}(X) := \lim_{K \to \infty} \frac{H(X^K)}{K} = \lim_{K \to \infty} \mathcal{L}_S^*\]

Idealmente si vorrebbe utilizzare blocchi di lunghezza sempre maggiore, solo che aumentando \(K\) aumenta anche la complessità
computazionale della codifica e lo spazio necessario per memorizzare le codewords associate a ciascun blocco. Nella pratica
serve scegliere il giusto compromesso tra lunghezza del blocco \(K\) e complessità computazionale.

\subsection{Codifica di Huffman}
\subsubsection*{Processo di codifica}
La codifica di Huffman è un algoritmo di compressione lossless che permette di ottenere uno dei codici ottimi per
codificare i simboli emessi da una sorgente discreta. L'algoritmo costruisce un albero binario in cui i simboli
compongono le foglie e le codewords sono ottenute unendo le etichette 0 e 1 degli archi che compongono i cammini
dalla radice alle foglie.

\subsubsection*{Criticità della codifica per simboli}
La codifica di Huffman presenta due principali criticità:
\begin{itemize}
	\item la lunghezza minima delle codewords è 1, quindi la lunghezza media del codice non può essere inferiore a 1,
	di conseguenza se l'entropia della sorgente è prossima allo zero (segnali molto sparsi e prevedibili), la codifica
	di Huffman non è ottimale
	\item la codifica di Huffman si basa solo sulla distribuzione di probabilità dei simboli, ma non tiene conto di
	possibili correlazioni tra simboli consecutivi emessi dalla sorgente, quindi non è ottimale per sorgenti in cui
	si hanno correlazioni tra simboli consecutivi (ad esempio pixel adiacenti di un'immagine)
\end{itemize}

\noindent
Per risolvere tali criticità si possono utilizzare le codifiche a blocchi, al costo di un aumento della complessità
esponenziale in funzione di \(K\). Per questo si utilizzano codifiche alternative per avvicinarsi al limite teorico
del tasso entropico della sorgente, come la codifica aritmetica o la codifica basata su dizionari.

\newpage


\subsection{Codifica aritmetica}
\subsubsection*{Processo di codifica per simboli}
La codifica aritmetica consiste nel suddividere l'intervallo \([0,1)\) in sottointervalli proporzionali alle probabilità
dei vari simboli emessi dalla sorgente. La codeword associata ad ogni simbolo è ottenuta convertendo in binario il numero
reale che rappresenta il punto medio del sottointervallo corrispondente al simbolo.

\subsubsection*{Estensione per codifica a blocchi}
Tale sistema risulta molto vantaggioso per codifiche a blocchi, in quanto per estendere i blocchi di simboli consecutivi
basta suddividere ricorsivamente il sottointervallo corrispondente al primo simbolo della sequenza con la distribuzione di
probabilità dei possibili simboli successivi. Tale distribuzione di probabilità può essere uguale a quella dei simboli presi
singolarmente (utile nel caso di simboli indipendenti) oppure è possibile usare le probabilità condizionate per avvicinarsi
al limite teorico del tasso entropico della sorgente.

\subsubsection*{Codifica adattiva}
A differenza di Huffman che richiede di conoscere a priori la distribuzione di probabilità dei simboli, la codifica aritmetica
permette di aggiornare dinamicamente il codice modificando le distribuzioni di probabilità e riadattando i sottointervalli
in base alla frequenza con cui i simboli vengono emessi.

\subsubsection*{Analisi delle prestazioni}
Il codice aritmetico per simboli singoli non è ottimale in quanto il limite superiore è più largo rispetto al codice di
Huffman, come indicato di seguito:
\[\mathcal{L}_A < H(X) + 2\]
Applicando la codifica aritmetica a blocchi di \(K\) simboli consecutivi, però, risulta più vantaggioso rispetto al codice
di Huffman in quanto la complessità aumenta linearmente con \(K\), e la lunghezza media per simbolo converge sempre al
tasso entropico della sorgente, come indicato di seguito:
\[\mathcal{L}_{A,S}^K < \frac{H(X^K)}{K} + \frac{2}{K} \quad \text{con } \mathcal{L}_{A,S}^K = \frac{\mathcal{L}_A}{K} \qquad\qquad \lim_{K \to \infty} \mathcal{L}_{A,S}^K = \lim_{K \to \infty} \frac{H(X^K)}{K} = \mathcal{H}(X)\]
Lo svantaggio principale è che la codifica aritmetica ha una complessità computazionale maggiore in quanto richiede di
svolgere operazioni tra numeri floating point con precisione sempre maggiore all'aumentare di \(K\) e necessita di tecnologie
proprietarie.

\subsection{Codifiche basate su dizionari - Lempel-Ziv, DEFLATE e ASN}
\subsubsection*{Funzionamento della codifica basata su dizionari (LZ77, LZ78, LZW)}
La codifica basata su dizionari sviluppata dagli algoritmi Lempel-Ziv (LZ77, LZ78) e Lempel-Ziv-Welch (LZW) si basa sul
riconoscimento di pattern e sequenze di simboli ripetute. In pratica le sequenze di simboli ricorrenti vengono sostituite
con un puntatore ad un dizionario o ad un punto precedente del messaggio stesso in cui è presente tale stringa.

\subsubsection*{Algoritmo DEFLATE}
L'algoritmo DEFLATE, utilizzato nelle compressioni ZIP, GZIP e PNG, combina la codifica basata su dizionari di LZ77 con
la codifica di Huffman per comprimere i puntatori delle stringhe ricorrenti. Lo svantaggio principale è che eredita tutti
i limiti della codifica di Huffman (non ottimalità per entropie basse e per simboli correlati).
Teoricamente sarebbe possibile sostituire la codifica di Huffman con la codifica aritmetica, ma si avrebbero notevoli
costi computazionali e di licenze.

\subsubsection*{Algoritmo ANS}
L'algoritmo ANS (Asymmetric Numeral Systems) si basa sul principio di funzionamento della codifica aritmetica, ma al posto
di suddividere l'intervallo \([0,1)\) in sottointervalli sempre più piccoli, si lavora con un insieme di numeri interi che
man mano crescono in grandezza ogni volta che si espande la dimensione del blocco di simboli da codificare secondo la
legge \(X' \approx X / P(s)\). In questo modo si evita di dover lavorare con numeri floating point ad altissima precisione.
Inoltre Jaroslaw Duda (l'ideatore dell'algoritmo ANS) ha rilasciato l'algoritmo in pubblico dominio senza brevetti contribuendo
ad una rapida diffusione e integrazione in molti software di compressione.

Esistono due varianti dell'algoritmo ANS:
\begin{itemize}
	\item \textbf{rANS (range ANS)}: le operazioni di aggiornamento degli stati interni vengono eseguite attraverso
	operazioni matematiche di moltiplicazioni e divisioni intere, risulta altamente parallelizzabile e permette di
	elaborare flussi di dati multipli; è usato in JPEG-XL e in Unreal Engine
	\item \textbf{tANS (table ANS)}: le operazioni di aggiornamento degli stati interni vengono eseguite attraverso
	una matrice di transizione pre-calcolata, risulta più veloce di rANS in quanto non ci sono operazioni matematiche,
	ma risulta meno parallelizzabile; è usato in Zstd di Meta, LZFSE di Apple
\end{itemize}

\subsection{Codifica Exp-Golomb}
\subsubsection*{Struttura base}
La codifica Exp-Golomb consiste nel codificare numeri interi con parole di codice di lunghezza crescente in funzione
del modulo del numero da codificare. Siccome non tiene conto della distribuzione di probabilità dei simboli, non è ottimale,
ma è molto semplice da implementare e risulta molto efficace per sorgenti in cui i simboli più probabili sono quelli
a valori bassi (es. per dati sparsi come i coefficienti DCT di un'immagine o di un video).

\subsubsection*{Exp-Golomb applicata a numeri naturali - Unsigned EGC}
Per codificare un numero naturale \(n\) con la codifica Exp-Golomb, si rappresenta il numero \(n+1\) in binario utilizzando
\(b = \lfloor \log_2(n+1) \rfloor + 1\) bit e si aggiunge un prefisso formato da una sequenza di 0 (leading zeros) di
lunghezza \(b-1\). La codifica di \(n = 0\) è 1. Siccome la rappresentazione binaria di \(n+1\) inizia sempre con
1, la sequenza di leading zeros non può venire confusa con la restante parte della codeword, inoltre suggerisce il
valore \(b\) facilitando la decodifica della codeword.

Ad esempio per il numero \(n = 8\), la codifica binaria di \(n+1=9\) è \(1001\), la lunghezza della sequenza di leading
zeros è \(b-1 = 3\), per cui la codeword finale è \(0001001\).

\subsubsection*{Exp-Golomb applicata a numeri interi - Signed EGC}
Per codificare un numero intero \(n\) con la codifica Exp-Golomb, si utilizza una mappatura \(m(n)\) dei numeri interi sui
numeri naturali per riutilizzare la codifica unsigned EGC.
\[m(n) = \begin{cases}
	2n -1 & \text{se } n > 0 \\
	-2n & \text{se } n \leq 0
\end{cases}\]

Ad esempio per il numero \(n = 8\), la mappatura è \(m(8) = 15\), la codifica binaria di \(m(8)+1=16\) è \(10000\), la
lunghezza della sequenza di leading zeros è \(b-1 = 4\), per cui la codeword finale è \(000010000\).

\subsection{Codifica per categoria e ampiezza}
La codifica per categoria e ampiezza è simile a Exp-Golomb, ma ha lunghezze di codeword più corte per numeri elevati al
costo di codeword più lunghe per numeri bassi.

Per codificare un numero intero \(n\) si utilizza un prefisso dato dal codice associato alla categoria \(k\) del numero
calcolata come \(k = \lceil \log_2 (|n|+1) \rceil\) e un suffisso dato dalla rappresentazione binaria di \(|n|\) con \(k\)
bit eventualmente in complemento a uno se \(n < 0\) (negazione bit-a-bit).

Ad esempio per il numero \(n = 8\), la categoria è \(k = 4\) associata al codice \(101\), l'ampiezza è \(1000\), per cui
la codeword finale è \(1011000\).

\newpage

\subsection{Codifica lossless predittiva}
\subsubsection*{Processo di codifica}
La codifica lossless predittiva consiste nell'utilizzare un predittore lineare per stimare il valore dei vari simboli
e salvare/trasmettere solo i residui di predizione codificati con un codificatore lossless. Essendo i residui meno
sparsi dei simboli originali, la codifica lossless dei residui risulta più efficiente (ha entropia più bassa) rispetto
alla codifica diretta dei simboli originali. Inoltre il residuo è un segnale sparso con valori più probabili attorno
allo zero, per cui la codifica Exp-Golomb risulta estremamente efficace e comparabile con la codifica ottima di Huffman.

I parametri utilizzati dal predittore lineare possono essere costanti per tutto il file o flusso, possono essere aggiornati
dinamicamente in base ai residui calcolati o possono essere calcolati per minimizzare l'errore quadratico medio dei residui
dei blocchi di simboli successivi.

\subsubsection*{Predittore semplice 1D}
Il predittore semplice assume che il simbolo corrente \(x_i\) sia uguale al simbolo precedente \(x_{i-1}\). Risulta
vantaggio per segnali semi-stazionari o con variazioni lente. Nelle immagini equivale ad una predizione lungo la direzione
orizzontale, ovvero risulta efficace per predire strutture orizzontali, ma non è efficace per predire strutture verticali
o diagonali.
\[\hat{X}(n) = X(n-1)\]

\subsubsection*{Predittore avanzato 2D}
Il predittore 2D, applicato alle immagini, assume che il pixel corrente \(x_{i,j}\) sia uguale al pixel adiacente a sinistra
\(x_{i,j-1}\) o al pixel adiacente sopra \(x_{i-1,j}\) in base a quale coppia di pixel adiacenti ha il valore più simile al
pixel lungo la diagonale \(x_{i-1,j-1}\). È capace di predire sia strutture orizzontali che verticali.
\[\hat{X}(i,j) = \begin{cases}
	x_{i,j-1} & \text{se } |x_{i,j} - x_{i,j-1}| \leq |x_{i,j} - x_{i-1,j}| \\
	x_{i-1,j} & \text{altrimenti}
\end{cases}\]

\subsubsection*{Applicazioni}
Tale codifica è molto efficace per sorgenti in cui i simboli consecutivi sono fortemente correlati ad esempio per
pixel adiacenti di un'immagine o per fotogrammi di un video.

Ad esempio un'immagine in scala di grigi con pixel a 8 bit, generalmente ha un'entropia di intorno ai 7 bpp (bit per
pixel) in quanto i livelli di grigio sono distribuiti abbastanza uniformemente. Un predittore 1D riesce a ridurre
l'entropia del residuo a circa 3.3 bpp e con un predittore 2D si riesce a ridurre ulteriormente fino a  circa 2.8 bpp.

\subsection{Codifica lossless tramite reti neurali}
\subsubsection*{Processo di codifica}
La codifica lossless tramite reti neurali si basa sull'addestramento di reti neurali (RNN, CNN, Transformers) come
predittori dei simboli emessi dalla sorgente. In questo modo si riescono a catturare correlazioni più complesse su
contesti di simboli più ampi e ottenere una distribuzione di probabilità altamente precisa per ogni blocco di simboli.
Tale distribuzione viene poi codificata tramite codifica entropica (aritmetica o ANS) e salvata come residuo di predizione.

\subsubsection*{Vantaggi e svantaggi}
Attraverso la predizione tramite reti neurali si riescono a raggiungere tassi di compressione molto vicini al limite
teorico del tasso entropico della sorgente, inoltre si riescono a catturare dipendenze non lineari e pattern a lungo
raggio.
Gli svantaggi principali sono l'elevata complessità computazionale data dalle operazioni in floating point per
l'inferenza della rete neurale che richiedono hardware specializzato (GPU, TPU) e tempi di compressione e decompressione
più lunghi rispetto alle codifiche tradizionali.
